% IEEE-formatted preliminary article (conference)
\documentclass[conference]{IEEEtran}
\usepackage[utf8]{inputenc}
\usepackage[T1]{fontenc}
\usepackage[spanish]{babel}
\usepackage{csquotes}
\usepackage{url}
\usepackage{graphicx}
\usepackage{cite}
\usepackage{hyperref}

\title{Desarrollo de una celda reconfigurable de CGRA}
\author{\IEEEauthorblockN{Duarte Sanchez David, Autor B, Autor C, Autor D, Autor E}
\IEEEauthorblockA{Grupo de 5\\
\{davidduarte28, autorB, autorC, autorD, autorE\}@estudiantec.cr}
}
\date{}

\begin{document}
\maketitle

\begin{abstract}
-Este documento es un borrador preliminar en formato IEEE que incluye título, autores, introducción y estado del arte. Aquí debe ir un resumen de 150--250 palabras describiendo objetivos, metodología y resultados esperados.
\end{abstract}

\section{Introducción}
La introducción debe motivar el problema, explicar su importancia y definir los objetivos del trabajo. Incluya contexto, preguntas de investigación y contribuciones principales.

\section{Antecedentes y Trabajo Relacionado}
Revisión de trabajos previos y estado del arte. Compare enfoques relevantes y destaque limitaciones que su trabajo aborda.

Cite referencias clave usando \cite{smith2019example,doe2020sample}.

\section{Conclusión provisional}
Resumen breve y próximos pasos.

\bibliographystyle{IEEEtran}
\bibliography{references}

\end{document}
